\documentclass[11pt]{article}

\usepackage[a4paper,margin=27mm]{geometry}
\usepackage[T1]{fontenc}
\usepackage[utf8]{inputenc}
\usepackage{lmodern}
\usepackage{microtype}
\usepackage{amsmath,amssymb,amsthm,bm}
\usepackage{booktabs}
\usepackage[hidelinks]{hyperref}

\hypersetup{
  pdftitle={A sharp weak-driving phase bound for symmetric complete Markov networks},
  pdfauthor={Xiaojian Zeng}
}

\newtheorem{theorem}{Theorem}
\newtheorem{proposition}[theorem]{Proposition}
\newtheorem{lemma}[theorem]{Lemma}
\newtheorem{corollary}[theorem]{Corollary}
\newcommand{\one}{\bm{1}}
\newcommand{\diag}{\operatorname{diag}}
\newcommand{\avg}[1]{\left\langle #1\right\rangle_T}

\title{A sharp weak-driving phase bound for symmetric complete Markov networks}
\author{Xiaojian Zeng\\
\small PBC School of Finance, Tsinghua University, Beijing 100190, China\\
\small \texttt{xiaojianzeng@foxmail.com}}
\date{}

\begin{document}
\maketitle

\begin{abstract}
Periodically driven Markov networks generally respond with a phase lag.  For a
complete $N$-state network with a common attempt rate and symmetric splitting of
energy differences, we prove a sharp bound on the leading weak-driving response.
For every $N\geq3$, every finite vector of energy offsets, and every finite real
Fourier drive, the quadratic response coefficients obey
$A_2\geq 2k_0B_2$.  The proof first expresses $A_2/B_2$ as a positive weighted
average of the nonzero relaxation rates.  The symmetrized generator is then reduced
to the rank-one matrix $\diag(1/q_i)-\one\one^{\mathsf T}$, whose smallest positive
eigenvalue is shown to be strictly larger than $2$ for every positive probability
vector $q$ when $N\geq3$.  We also classify equality: at finite offsets it occurs
only for zero-response drives, whereas every drive that changes an energy difference
gives a strict inequality.  An explicit two-heavy-state family approaches the bound
as the remaining stationary weights vanish, proving that the constant $2$ is optimal
but unattained by a response-active finite-offset system.  The result is confined to
the weak-amplitude, equal-prefactor, symmetric-splitting complete graph; broader
finite-amplitude and heterogeneous-network bounds remain open.
\end{abstract}

\noindent\textbf{Keywords:} stochastic thermodynamics; periodic driving; phase lag;
Markov networks; spectral gap; rank-one perturbation

\section{Introduction}

A periodic modulation of a stochastic system is followed by a delayed response.
For finite-state Markov dynamics, Deg\"unther, Koyuk and Seifert introduced a
generalized phase that remains meaningful for multi-harmonic driving and related it
to cycle-averaged contractions of the probability distribution with the first and
second time derivatives of the energies \cite{Deguenther2022}.  They established a
tight two-state bound.  For multi-state networks with a common rate amplitude they
reported numerical evidence and left the corresponding general proof open; their
heterogeneous-amplitude statement is broader still.

Here we resolve a sharply delimited part of that problem.  We consider the complete
graph $K_N$, a common attempt rate $k_0$, symmetric splitting, finite energy offsets,
and the leading nonvanishing order in the driving amplitude.  The response problem
then becomes a spectral-gap problem for a nonuniform diagonal matrix minus a rank-one
matrix.  Rank-one perturbations are a standard tool in graph spectral theory
\cite{KleeStamps2020}, but the offset-dependent matrix here has kernel proportional to
a nonuniform vector, and immediate interlacing does not settle the regime in which a
single auxiliary weight exceeds one half.

Our main result is
\begin{equation}
  A_2\geq 2k_0B_2 .                                           \label{eq:main-bound}
\end{equation}
For $N\geq3$ and finite offsets, the inequality is strict whenever the drive changes
at least one energy difference.  Equality is precisely the zero-response class.  The
constant is nevertheless sharp because a family with two equally weighted dominant
states converges to the symmetric two-state mechanism.  Besides proving
\eqref{eq:main-bound}, this identifies why adding finitely occupied states cannot
increase the weak-drive lag beyond that limiting mechanism.

Section~\ref{sec:model} defines the physical domain and derives the modal response.
Section~\ref{sec:gap} proves the exact spectral inequality.  Section~\ref{sec:phase}
classifies equality and establishes sharpness.  We close with the physical scope and
the unresolved extensions.

\section{Model and weak response}                         \label{sec:model}

\subsection{Complete detailed-balance network}

Let $N\geq3$, $k_0>0$, and $T>0$, with fundamental frequency
$\Omega=2\pi/T$.  Energies are measured in units in which
$k_{\mathrm B}T_{\mathrm bath}=1$.  The rate from state $i$ to state $j$ is
\begin{equation}
 k_{ij}^{(\varepsilon)}(t)
 =k_0\exp\!\left[\frac{E_i^{(\varepsilon)}(t)
                         -E_j^{(\varepsilon)}(t)}{2}\right],
 \qquad
 E_i^{(\varepsilon)}(t)=E_i^0+\varepsilon e_i(t),          \label{eq:rates}
\end{equation}
where all $E_i^0$ are finite and each $e_i$ is a finite real Fourier series of
period $T$.  With $p(t)$ a normalized column vector, the master equation is
\begin{equation}
 \dot p_i(t)=\sum_{j\ne i}
 \left[p_j(t)k_{ji}(t)-p_i(t)k_{ij}(t)\right].              \label{eq:master}
\end{equation}
All rates are positive, so the long-time $T$-periodic state is unique.

For that periodic state, define
\begin{equation}
 \mathcal A(\varepsilon)=\frac1T\int_0^T
 p^{(\varepsilon)}(t)\mathbin{\cdot}\ddot E^{(\varepsilon)}(t)\,dt,
 \qquad
 \mathcal B(\varepsilon)=\frac1T\int_0^T
 p^{(\varepsilon)}(t)\mathbin{\cdot}\dot E^{(\varepsilon)}(t)\,dt. \label{eq:AB}
\end{equation}
The constant and linear terms vanish by periodicity.  We use the coefficient
convention
\begin{equation}
 A_2=\lim_{\varepsilon\to0}\frac{\mathcal A(\varepsilon)}{\varepsilon^2},
 \qquad
 B_2=\lim_{\varepsilon\to0}\frac{\mathcal B(\varepsilon)}{\varepsilon^2}. \label{eq:A2B2}
\end{equation}
Thus $A_2$ has units of inverse time squared, $B_2$ has units of inverse time,
and $A_2/B_2$ is a relaxation rate when $B_2>0$.

At $\varepsilon=0$, introduce
\begin{equation}
 x_i=e^{-E_i^0/2},\qquad S=\sum_i x_i,\qquad
 Q=\sum_i x_i^2,\qquad \pi_i=\frac{x_i^2}{Q}.               \label{eq:weights}
\end{equation}
The equilibrium vector is $\pi$.  If $L_0$ denotes the column-conservative
generator in \eqref{eq:master} and $\Pi=\diag(\pi_i)$, detailed balance gives the
real symmetric, symmetrized negative generator
\begin{equation}
 M=-\Pi^{-1/2}L_0\Pi^{1/2}
  =k_0\left[S\diag(x_1^{-1},\ldots,x_N^{-1})
              -\one\one^{\mathsf T}\right],
 \qquad Mx=0.                                                \label{eq:M}
\end{equation}

\subsection{Spectral representation of the response}

Write the first-order periodic response as
$p^{(\varepsilon)}=\pi+\varepsilon r+O(\varepsilon^2)$ and set
$r=\Pi^{1/2}z$.  For
\begin{equation}
 \widetilde e(t)=e(t)-\one\,\pi^{\mathsf T}e(t),
 \qquad h(t)=\Pi^{1/2}\widetilde e(t),                       \label{eq:centered}
\end{equation}
The tangent to the instantaneous equilibrium distribution is
$\delta\pi=-\Pi\widetilde e$.  Differentiating
$L(E)\pi(E)=0$ therefore gives
$L'[e]\pi=L_0\Pi\widetilde e$.  Substitution in the first-order master equation
$\dot r=L_0r+L'[e]\pi$, followed by the similarity transform in \eqref{eq:M}, gives
\begin{equation}
 \dot z=-Mz-Mh .                                             \label{eq:z}
\end{equation}
Normalization makes $z$ orthogonal to $\ker M$, so only the positive modes enter.

To make all Fourier factors explicit, write
\begin{equation}
 e(t)=e^{(0)}+\sum_{n=1}^{H}\operatorname{Re}
       \left[c_n e^{\mathrm i\omega_nt}\right],
 \qquad \omega_n=n\Omega,                                   \label{eq:fourier}
\end{equation}
and let $h_n=\Pi^{1/2}(c_n-\one\pi^{\mathsf T}c_n)$.  Choose an
orthonormal positive-mode eigenbasis $Mu_\mu=\lambda_\mu u_\mu$, and put
$a_{n\mu}=u_\mu^{\mathsf T}h_n$.

\begin{proposition}[Modal response]
For the model \eqref{eq:rates}--\eqref{eq:fourier},
\begin{align}
 A_2&=\frac12\sum_{n,\mu>0}
 \frac{\omega_n^2\lambda_\mu^2|a_{n\mu}|^2}
      {\lambda_\mu^2+\omega_n^2},                           \label{eq:Amodal}\\
 B_2&=\frac12\sum_{n,\mu>0}
 \frac{\omega_n^2\lambda_\mu|a_{n\mu}|^2}
      {\lambda_\mu^2+\omega_n^2}.                           \label{eq:Bmodal}
\end{align}
Consequently, when $B_2>0$, $A_2/B_2$ is a positive weighted average of the
nonzero eigenvalues of $M$.
\end{proposition}

\begin{proof}
For a mode with frequency $\omega_n$ and eigenvalue $\lambda_\mu$,
\eqref{eq:z} gives the complex response amplitude
\begin{equation}
 Z_{n\mu}=-\frac{\lambda_\mu}{\lambda_\mu+\mathrm i\omega_n}
 a_{n\mu}.                                                   \label{eq:Z}
\end{equation}
Expanding \eqref{eq:AB} to order $\varepsilon^2$ yields
$A_2=\avg{r\cdot\ddot e}$ and $B_2=\avg{r\cdot\dot e}$.
Because $z$ is orthogonal to $\Pi^{1/2}\one$, these contractions equal
$\avg{z\cdot\ddot h}$ and $\avg{z\cdot\dot h}$.  Taking the real-period
average of each complex harmonic in \eqref{eq:Z} gives
\eqref{eq:Amodal} and \eqref{eq:Bmodal}.  Every weight in the latter is
nonnegative and is positive exactly when the corresponding $a_{n\mu}$ is nonzero.
\end{proof}

For $N=2$ at equal offsets, the sole positive eigenvalue is $2k_0$, recovering
$A_2=2k_0B_2$ and the symmetric two-state phase reported in
Ref.~\cite{Deguenther2022}.  The remaining task is therefore an all-$N$ lower bound
on the positive spectrum of \eqref{eq:M}.

\section{A sharp spectral gap}                              \label{sec:gap}

Normalize $x$ in the one-norm by defining $q_i=x_i/S$.  Then $q_i>0$,
$\sum_iq_i=1$, and
\begin{equation}
 \frac{M}{k_0}=A(q),\qquad
 A(q)=\diag(q_1^{-1},\ldots,q_N^{-1})-\one\one^{\mathsf T}. \label{eq:Aq}
\end{equation}

\begin{lemma}[Positivity and kernel]
$A(q)$ is positive semidefinite and $\ker A(q)=\operatorname{span}\{q\}$.
\end{lemma}

\begin{proof}
For $y\in\mathbb R^N$, put $b_i=y_i/q_i$.  Then
\begin{align}
 y^{\mathsf T}A(q)y
 &=\sum_iq_i b_i^2-\left(\sum_iq_i b_i\right)^2 \\
 &=\sum_{i<j}q_iq_j(b_i-b_j)^2\geq0.                         \label{eq:variance}
\end{align}
Since every $q_i$ is positive, equality holds exactly when all $b_i$ are equal,
or equivalently when $y$ is proportional to $q$.
\end{proof}

\begin{theorem}[Strict finite-offset gap]
\label{thm:gap}
For every $N\geq3$ and every positive vector $q$ with $\sum_iq_i=1$,
\begin{equation}
 \lambda_{\min}^{+}\!\left(A(q)\right)>2.                  \label{eq:gap}
\end{equation}
Equivalently, every nonzero relaxation rate in \eqref{eq:M} is larger than
$2k_0$.
\end{theorem}

\begin{proof}
Relabel the coordinates so that $q_1\geq q_2\geq\cdots\geq q_N$, and set
$d_i=1/q_i$, so $d_1\leq\cdots\leq d_N$.  Let the eigenvalues of $A(q)$ be
$0=\lambda_1<\lambda_2\leq\cdots\leq\lambda_N$.  Rank-one interlacing gives
\begin{equation}
 0=\lambda_1\leq d_1\leq\lambda_2\leq d_2\leq\cdots .      \label{eq:interlace}
\end{equation}

If $q_1<1/2$, then $d_1>2$, and \eqref{eq:interlace} immediately gives
$\lambda_2>2$.

If $q_1=1/2$, interlacing gives $\lambda_2\geq2$.  The value $2$ cannot be an
eigenvalue.  Indeed, $A(q)y=2y$ would imply
\begin{equation}
 (d_i-2)y_i=c,\qquad c=\one^{\mathsf T}y.                    \label{eq:exclude2}
\end{equation}
The first coordinate gives $c=0$.  Since $N\geq3$ and all $q_i$ are positive,
$q_i<1/2$ for $i\geq2$, so \eqref{eq:exclude2} gives $y_i=0$ for those indices.
It then also gives $y_1=c=0$, excluding a nonzero eigenvector.

It remains to consider $q_1>1/2$.  This largest coordinate is unique.  Put
$s=1-q_1\in(0,1/2)$; then $d_1<2<d_2$.  On $(d_1,d_2)$ the determinant lemma
gives the secular function
\begin{equation}
 F(\lambda)=1-\sum_{i=1}^N\frac{1}{d_i-\lambda},
 \qquad F'(\lambda)=-\sum_{i=1}^N\frac{1}{(d_i-\lambda)^2}<0. \label{eq:secular}
\end{equation}
It decreases from $+\infty$ to $-\infty$, and its unique zero is $\lambda_2$.

For $g(t)=t/(1-2t)$ on $[0,1/2)$, direct algebra gives
\begin{equation}
 g(a+b)-g(a)-g(b)
 =\frac{4ab(1-a-b)}
 {(1-2a)(1-2b)(1-2a-2b)}>0                                 \label{eq:superadd}
\end{equation}
whenever $a,b>0$ and $a+b<1/2$.  The $N-1\geq2$ positive numbers
$q_2,\ldots,q_N$ sum to $s$, so repeated strict superadditivity yields
$\sum_{i\geq2}g(q_i)<g(s)$.  Hence
\begin{align}
 \sum_{i=1}^N\frac{1}{d_i-2}
 &=-\frac{1-s}{1-2s}+\sum_{i\geq2}g(q_i) \\
 &<-\frac{1-s}{1-2s}+\frac{s}{1-2s}=-1.                    \label{eq:F2}
\end{align}
Thus $F(2)>0$.  Since $F$ is strictly decreasing and vanishes at $\lambda_2$,
we conclude that $\lambda_2>2$.  The three cases exhaust the positive simplex.
\end{proof}

The strictness in \eqref{eq:gap} depends on both positivity of every $q_i$ and
$N\geq3$.  It is therefore a finite-offset statement; it does not assert a uniform
positive margin above $2$ as the boundary of the simplex is approached.

\section{Phase bound, equality, and optimality}             \label{sec:phase}

\begin{corollary}[Weak-response phase bound]
For every finite-offset model \eqref{eq:rates} with $N\geq3$ and every finite real
Fourier drive,
\begin{equation}
 A_2-2k_0B_2
 =\frac12\sum_{n,\mu>0}
 \frac{\omega_n^2\lambda_\mu(\lambda_\mu-2k_0)
       |a_{n\mu}|^2}{\lambda_\mu^2+\omega_n^2}\geq0.        \label{eq:response-diff}
\end{equation}
If at least one time-dependent energy difference is nonconstant, then
$A_2>2k_0B_2>0$.  Equality holds exactly when all positive-frequency Fourier
coefficients are common to every state, in which case $A_2=B_2=0$.
\end{corollary}

\begin{proof}
Equation~\eqref{eq:response-diff} follows by subtracting twice $k_0$ times
\eqref{eq:Bmodal} from \eqref{eq:Amodal}.  Every nonzero factor
$\lambda_\mu-2k_0$ is strictly positive by Theorem~\ref{thm:gap}.  Moreover, all
$a_{n\mu}$ vanish if and only if
$c_n-\one\pi^{\mathsf T}c_n=0$ for every $n\geq1$, which is precisely a common
time-dependent energy shift.  Such a shift leaves every rate invariant.  Arbitrary
zero-frequency shifts likewise make no contribution to the cycle averages.
\end{proof}

In the phase convention of Ref.~\cite{Deguenther2022}, a nonzero weak response obeys
\begin{equation}
 \tan\Delta\varphi=\Omega\frac{B_2}{A_2}
 <\frac{\Omega}{2k_0}.                                      \label{eq:phase}
\end{equation}
Thus every response-active finite-offset $N\geq3$ system lies strictly below the
symmetric two-state ceiling.  The strict inequality does not mean that the constant
can be improved uniformly.

\begin{proposition}[Sharpness]
For each fixed $N\geq3$, the constant $2$ in \eqref{eq:main-bound} is the optimal
uniform constant over response-active systems with finite offsets, but it is not
attained.
\end{proposition}

\begin{proof}
For $0<\delta\leq1$, take
\begin{equation}
 x_1=x_2=1,\qquad x_3=\cdots=x_N=\delta,
 \qquad S=2+(N-2)\delta.                                    \label{eq:boundary-family}
\end{equation}
The vector $v=(1,-1,0,\ldots,0)^{\mathsf T}$ is orthogonal to $\one$ and is an
exact eigenvector of \eqref{eq:M} with eigenvalue $k_0S$.  The two smallest diagonal
entries of $S\diag(1/x_i)$ are both $S$; interlacing therefore shows that this is the
smallest positive eigenvalue.  A single real harmonic proportional to $v$ is centered
because $\pi_1=\pi_2$, and it excites only this mode.  Consequently,
\begin{equation}
 \frac{A_2}{B_2}=k_0\left[2+(N-2)\delta\right].              \label{eq:sharp-ratio}
\end{equation}
Every $\delta>0$ corresponds to finite energies.  As $\delta\downarrow0$, the
energies of states $3,\ldots,N$ diverge and \eqref{eq:sharp-ratio} approaches
$2k_0$ from above.  Hence any larger uniform constant fails for sufficiently small
positive $\delta$, whereas Theorem~\ref{thm:gap} excludes attainment at a finite offset.
\end{proof}

\section{Discussion and limitations}

The response formulas separate the physical and spectral parts of the result.  The
drive, frequency, and observable determine nonnegative modal weights, whereas the
offsets determine the relaxation spectrum.  Complete connectivity with a common
attempt rate forces every response-active mode above $2k_0$.  The sharp family shows
that the slowest allowed mechanism is not a genuinely multi-state interior point: it
is approached when only two equally weighted states remain dynamically relevant.

The nonuniform gap is not obtained from the ordinary unweighted complete-graph
spectrum alone.  When $\max_i q_i\leq1/2$, interlacing almost settles the question;
when one weight is dominant, the sign of the first secular root is controlled by the
strict superadditivity identity \eqref{eq:superadd}.  This is the only regime in which
the elementary diagonal lower bound is insufficient.

Several restrictions are essential.  First, the statement concerns only the
coefficients defined by the limit $\varepsilon\to0$; it is not a finite-amplitude
theorem.  Second, all edges have the same prefactor $k_0$ and use symmetric splitting.
Third, the graph is complete and all offsets are finite.  Missing links,
edge-dependent prefactors, asymmetric splitting, transient states, and singular
offsets are outside the proof.  In particular, the broader multi-state conjectures in
Ref.~\cite{Deguenther2022} remain unresolved.

\section{Conclusion}

We have proved the weak-driving inequality $A_2\geq2k_0B_2$ for every finite-offset
symmetric complete $N$-state Markov network with $N\geq3$.  The exact spectral
reduction yields more than the non-strict bound: a response-active interior system is
always strict, equality is exactly the zero-response gauge class, and the constant is
an unattained sharp infimum.  The result turns the two-state phase scale into a precise
limiting benchmark for this complete multi-state model while leaving finite-amplitude
and heterogeneous-network extensions as distinct open problems.

\section*{Data and code availability}

No experimental data were generated.  The analytic result is supported by accompanying
Python programs that (i) reconstruct the weak-response formula in both the probability
and detailed-balance bases, (ii) perform exact-rational finite-instance gap checks, and
(iii) verify the secular and superadditivity identities symbolically.  These programs
are verification aids and do not replace the proof.  A repository location and
persistent identifier should be added by the author before submission.

\section*{Acknowledgments and generative-AI disclosure}

OpenAI Codex using model \texttt{gpt-5.6-sol} assisted with research ideation, algebraic calculations,
verification code, literature discovery, and manuscript drafting.  All scientific
claims, references, calculations, and wording require the named author's independent
review and approval before submission.  No AI system is listed as an author.

\section*{Conflict of interest}

The author must confirm the conflict-of-interest declaration before submission.

\bibliographystyle{unsrt}
\bibliography{references}

\end{document}
